% ============================================================
% SST CANON v0.8.15 PATCH SET
% Topic: triadic channel locking, no-monopole transport-stress audit,
%        and SST-73 alpha_grav dimensional correction.
% Prepared for insertion/audit; not an automatic overwrite.
% ============================================================

% ------------------------------------------------------------
% PATCH A: Insert into SST_CANON-v0.8.15.tex
% Location: subsection Triadic Gravity-Response Corollary,
% after the sentence:
% "The three modes may share a common source state but are not mutually identical."
% ------------------------------------------------------------

\paragraph{Local Euler locking of pressure and optical projections.}
Under the additional assumptions of a passive, stationary, inviscid Euler
region with constant \(\rhoF\), the pressure and optical projections are not
independent.  Bernoulli reconstruction gives
\begin{align}
    \delta p_{\mathrm{swirl}}
    =
    -\frac{1}{2}\rhoF\lVert\uswirl\rVert^2,
\end{align}
while the clock/optical closure gives
\begin{align}
    n_\gamma
    =
    \left(
        1-rac{\lVert\uswirl\rVert^2}{c^2}
    \right)^{-1/2}
    =
    1+rac{\lVert\uswirl\rVert^2}{2c^2}
    +O\!\left(\frac{\lVert\uswirl\rVert^4}{c^4}\right).
\end{align}
Eliminating \(\lVert\uswirl\rVert^2\) yields the parameter-free leading-order
locking relation
\begin{align}
    \boxed{
    \delta n_\gamma
    =
    -\frac{\delta p_{\mathrm{swirl}}}{\rhoF c^2}
    }
    \label{eq:canonical_pressure_optical_locking}
\end{align}
with validity restricted to the passive Euler regime.  Forcing, damping,
viscosity, thermal gradients, acoustic drive, plasma response, or explicit
Mode-I terms such as \(\Gamma_{\swirlarrow}\mathbf{u}_{\mathrm{plume}}\)
break the assumptions and move the prediction to the research-track transfer
model.  Thus the refined diagnostic statement is not three fully independent
channels, but two locally locked projections plus one nonlocal Poisson/bulk
response whenever Eq.~\eqref{eq:canonical_weak_gravity_closure} is invoked.

With \(\rhoF=7.0\times10^{-7}\,\mathrm{kg\,m^{-3}}\),
\begin{align}
    \frac{1}{\rhoF c^2}
    &=
    1.58950008\times10^{-11}\,\mathrm{Pa^{-1}},
    \\
    \frac{1}{2}\rhoF\vchar^2
    &=
    4.1877439\times10^5\,\mathrm{Pa},
    \\
    \delta n_\gamma
    &=
    6.6565\times10^{-6}
\end{align}
at canonical swirl speed, consistent with Eq.~\eqref{eq:canonical_emg_numerical_hierarchy}.

% ------------------------------------------------------------
% PATCH B: Insert into SST_CANON-v0.8.15-research-track.tex
% Location: after section Flame-Lock, Photonless-Caustic, and Elastic-Shell
% Diagnostics, before section Atomic Gravity Closure and Condensate Coherence.
% ------------------------------------------------------------

\section{Research Track: No-Monopole Audit of the Transport-Stress Gravity Candidate}
\label{sec:rt_no_monopole_transport_stress_audit}

\paragraph{Status.}
\textbf{[RESEARCH TRACK / DERIVED CONSTRAINT].}
This section records a constraint on the SST-73 transport-stress gravity
candidate.  It does not modify the canon-level weak-field closure
\(\nabla^2\Phi_{\swirlarrow}=4\pi G_{\mathrm{swirl}}\rhoM\).  It shows that a
double-divergence stress source belongs to the local pressure/stress channel
unless an additional non-divergence density closure or boundary source is
specified.

\subsection{Pressure equivalence lemma}
For an incompressible inviscid medium with constant \(\rhoF\),
\begin{align}
    \partial_i u_i=0,
\end{align}
and no external body force, the Euler equation is
\begin{align}
    \rhoF\left(\partial_t u_i + u_j\partial_j u_i\right)
    =
    -\partial_i p .
\end{align}
Taking the divergence and using \(\partial_i u_i=0\) gives
\begin{align}
    \partial_i\partial_j\left(\rhoF u_i u_j\right)
    =
    -\nabla^2p .
    \label{eq:rt_pressure_equivalence_euler_identity}
\end{align}
Therefore the SST-73 transport-stress candidate
\begin{align}
    \nabla^2\Phi_{\mathrm{tr}}
    =
    \lambda\,\partial_i\partial_j\left(\rhoF u_i u_j\right)
\end{align}
is equivalent to
\begin{align}
    \nabla^2\left(\Phi_{\mathrm{tr}}+\lambda p\right)=0 .
\end{align}
For decaying exterior boundary conditions,
\begin{align}
    \boxed{
    \Phi_{\mathrm{tr}}
    =
    -\lambda\left(p-p_\infty\right)
    }
    \label{eq:rt_transport_pressure_equivalence}
\end{align}
up to an irrelevant harmonic constant.  Candidate C is therefore not an
independent long-range gravity mechanism in the smooth incompressible bulk; it
is the Euler pressure channel written in Poisson form.

\subsection{No-monopole theorem for smooth compact stress sources}
Let
\begin{align}
    S(\mathbf{x})
    =
    \partial_i\partial_j\Sigma_{ij},
    \qquad
    \Sigma_{ij}=\rhoF u_i u_j,
\end{align}
with \(\Sigma_{ij}\) smooth and compactly supported, or decaying sufficiently
fast at infinity.  The monopole moment is
\begin{align}
    Q_0
    =
    \int_{\mathbb{R}^3}S(\mathbf{x})\,dV
    =
    \oint_{\partial\mathbb{R}^3}\partial_j\Sigma_{ij}\,dA_i
    =0 .
\end{align}
The dipole moment also vanishes after integration by parts:
\begin{align}
    Q_k
    =
    \int_{\mathbb{R}^3}x_k S(\mathbf{x})\,dV
    =0 .
\end{align}
Hence the first possible nonzero exterior multipole is quadrupolar:
\begin{align}
    \Phi_{\mathrm{tr}}(r)=O(r^{-3}),
    \qquad
    \lVert\nabla\Phi_{\mathrm{tr}}\rVert=O(r^{-4}).
\end{align}
This cannot supply a Newtonian \(1/r^2\) acceleration by itself.

\subsection{Boundary and singular-core exception}
The no-monopole result assumes a smooth regularized field on all of
\(\mathbb{R}^3\).  If vortex tubes are excised, singular cores are retained,
or material boundaries/shells are present, the surface term
\begin{align}
    \oint_{\partial V}\partial_j\Sigma_{ij}\,dA_i
\end{align}
need not vanish.  Such terms are boundary-stress or shell-response physics
\(\mathcal{R}_{\Pi}\), not bulk Poisson gravity.  A nonzero Newtonian monopole
therefore requires either a genuine scalar density source, such as
\(\rhoM=\rhoE/c^2\), or an explicitly modeled boundary source.

\subsection{Equivalence-principle consequence of the density source}
If the gravitational source density is the same effective mass density that
defines inertial mass,
\begin{align}
    \rhoM^{\mathrm{eff}}
    =
    \frac{\rhoE}{c^2},
    \qquad
    \rhoE=\frac{1}{2}\rhoF\lVert\uswirl\rVert^2,
\end{align}
then
\begin{align}
    M_{\mathrm{grav}}
    =
    \int\rhoM^{\mathrm{eff}}\,dV
    =
    \frac{1}{c^2}\int\rhoE\,dV
    =
    M_{\mathrm{inert}} .
\end{align}
For full topological SST states this statement must be applied to the same
coarse-grained mass functional used for inertial mass, including topology,
geometric gate, and clock-impedance factors.  The numerical value of
\(G_{\mathrm{swirl}}\) remains \textbf{[CALIBRATED]} until derived
independently of orthodox \(G\) or \(t_p\).

\subsection{Minimal numerical tests}
A verification script for this sector must include:
\begin{enumerate}
    \item a smooth compact divergence-free test field for which
    \(\int S\,dV\approx0\) and \(\int x_kS\,dV\approx0\);
    \item a far-field multipole fit verifying
    \(\Phi_{\mathrm{tr}}\sim r^{-3}\), not \(r^{-1}\);
    \item a pressure-equivalence test verifying
    \(\Phi_{\mathrm{tr}}+\lambda p\) is harmonic/constant under the stated
    boundary conditions;
    \item a separate density-source Poisson test verifying
    \(\Phi\sim -G_{\mathrm{swirl}}M/r\) and
    \(\lVert\mathbf{g}\rVert\sim r^{-2}\) for compact positive
    \(\rhoM^{\mathrm{eff}}\).
\end{enumerate}

% ------------------------------------------------------------
% PATCH C: Insert into SST-73_A_Poisson-Type_Gravity_Program_from_Organized_Swirl_Transport.tex
% Location 1: after interpretation of Eq. (poisson-tensor), before elevating Candidate C.
% ------------------------------------------------------------

\paragraph{Critical note: pressure equivalence and no bulk monopole.}
For constant-density incompressible Euler flow without external body force,
\begin{align}
    \partial_i\partial_j\left(\rhoF u_i u_j\right)
    =
    -\nabla^2p .
\end{align}
Consequently the tensorial Candidate C is pressure-equivalent:
\begin{align}
    \nabla^2\Phi_{\grav}
    =
    \lambda\partial_i\partial_j(\rhoF u_i u_j)
    \quad\Longrightarrow\quad
    \nabla^2(\Phi_{\grav}+\lambda p)=0 .
\end{align}
With decaying exterior boundary conditions this gives
\begin{align}
    \Phi_{\grav}=-\lambda(p-p_\infty) .
\end{align}
For smooth compact stress fields the integrated source
\(\int\partial_i\partial_j\Sigma_{ij}\,dV\) vanishes by Gauss' theorem, and the
dipole moment also vanishes by integration by parts.  The first possible
far-field moment is quadrupolar.  Candidate C should therefore be retained as a
local pressure/stress diagnostic unless an explicit boundary source or a
non-divergence scalar density source is added.  This is the same mathematical
obstruction that appears in Lighthill's quadrupole source structure for
aerodynamic sound.

% Location 2: after Eq. (rho-g-axis-compact) or after Eq. (poisson-practical).

\paragraph{Dimensional note on \(\alpha_{\grav}\).}
In the practical axisymmetric proxy
\begin{align}
    \rho_{\grav}^{(\SST)}
    =
    \alpha_{\grav}\nabla^2(\rhoF u_\theta^2),
\end{align}
\(\alpha_{\grav}\) is not dimensionless.  Since
\begin{align}
    [\rho_{\grav}^{(\SST)}]=\mathrm{kg\,m^{-3}},
    \qquad
    [\nabla^2(\rhoF u_\theta^2)]
    =
    \mathrm{kg\,m^{-3}\,s^{-2}},
\end{align}
one must have
\begin{align}
    \boxed{[\alpha_{\grav}]=\mathrm{s^2}.}
\end{align}
Equivalently, in
\(\nabla^2\Phi_{\grav}=4\pi G_{\eff}\alpha_{\grav}\nabla^2(\rhoF u_\theta^2)\),
\(\alpha_{\grav}\) supplies the missing time-scale squared.  A natural but not
yet derived candidate is
\begin{align}
    \alpha_{\grav}
    \sim
    \left(\frac{\rc}{\vchar}\right)^2
    \approx
    1.6595\times10^{-42}\,\mathrm{s^2},
\end{align}
which is \textbf{[SPECULATIVE]} until obtained from a closure theorem or
calibration.

\begin{thebibliography}{99}
\bibitem{Lighthill1952}
M.~J. Lighthill (1952),
``On Sound Generated Aerodynamically. I. General Theory,''
\emph{Proceedings of the Royal Society of London. Series A, Mathematical and Physical Sciences}
\textbf{211}(1107), 564--587.
DOI: \href{https://doi.org/10.1098/rspa.1952.0060}{10.1098/rspa.1952.0060}.
\end{thebibliography}
